% Original analytic teaching diagram, 2026-11-12.
% Steps of the characteristic-function proof of the i.i.d. central limit theorem,
% each box tagged with the hypothesis it uses (independence, centring, variance, continuity theorem).
\definecolor{VB}{HTML}{7188A3}\definecolor{VBF}{HTML}{E5EAF0}
\definecolor{VG}{HTML}{829C8B}\definecolor{VGF}{HTML}{E8EEE9}
\definecolor{VY}{HTML}{B4A16F}\definecolor{VYF}{HTML}{F1ECDD}
\definecolor{VR}{HTML}{A7787D}\definecolor{VRF}{HTML}{F0E5E6}
\definecolor{VP}{HTML}{9588A5}\definecolor{VPF}{HTML}{ECE8F0}
\definecolor{VI}{HTML}{3B4252}\definecolor{VM}{HTML}{6B7280}\definecolor{VL}{HTML}{4C4D4F}
\begin{tikzpicture}[x=1mm,y=1mm,>=Stealth,font=\figurefont\fontsize{8}{10}\selectfont,
lab/.style={text=VM,font=\figurefont\fontsize{7.5}{9}\selectfont},
box/.style={draw=VB,fill=VBF,line width=.6pt,rounded corners=1.2mm,minimum height=10mm,minimum width=31mm,align=center,text=VI},
goal/.style={draw=VR,fill=VRF,line width=.8pt,rounded corners=1.2mm,minimum height=10mm,minimum width=31mm,align=center,text=VI},
arr/.style={->,draw=VL,line width=.6pt}]
\node[box] (b1) at (0,0) {$\phi_{Z_n}(t)=\phi_Y\bigl(\tfrac{t}{\sqrt{n}}\bigr)^{n}$};
\node[box] (b2) at (38,0) {$\phi_Y(s)=1-\tfrac{s^{2}}{2}+R(s)$};
\node[box] (b3) at (76,0) {$n\,\bigl\lvert R\bigl(\tfrac{t}{\sqrt{n}}\bigr)\bigr\rvert\to0$};
\node[box] (b4) at (19,-24) {$\phi_{Z_n}(t)\to\mathrm{e}^{-t^{2}/2}$};
\node[goal] (b5) at (57,-24) {$Z_n\xrightarrow{\mathrm{d}}\mathcal{N}(0,1)$};
\draw[arr] (b1) -- (b2);
\draw[arr] (b2) -- (b3);
\draw[arr] (b3.south) -- ++(0,-5) -| (b4.north);
\draw[arr] (b4) -- (b5);
\node[lab,below=1.2mm of b1] {独立：乘积};
\node[lab,below=1.2mm of b2] {中心化、方差为一};
\node[lab,below=1.2mm of b3,xshift=-9mm] {二阶矩有限};
\node[lab,below=1.2mm of b4] {逐点收敛};
\node[lab,below=1.2mm of b5] {连续性定理};
\end{tikzpicture}
